\begin{table}[htbp]\centering
\caption{Policy implications and application boundaries}\label{tab:key_findings}
\begin{minipage}{.95\textwidth}\fontsize{8}{10}\selectfont\setstretch{1}\setlength{\tabcolsep}{2pt}\renewcommand{\arraystretch}{1.08}
\par\noindent\textbf{Evidence for implementation}\par\smallskip
\begin{tabular*}{\linewidth}{@{\extracolsep{\fill}}>{\raggedright\arraybackslash}p{.32\linewidth}>{\raggedright\arraybackslash}p{.62\linewidth}@{}}\toprule
Implication & Evidence and scope \\
\midrule
Regulate fares together with service & At the central anchor, cap 1.10 cuts commercial welfare losses by 78.8/51.4/57.3\%, retaining 74.5/84.5/60.3\% of profit gains. \\
Compare participation floors with direct fare restraint & Strict HN/WC retention attains 69.2/94.9\% of certified reference welfare benchmarks. Cap 1.10 alone leaves welfare below P0; cap 1.00 restores positive welfare. Joint caps and floors give corridor-specific profit--welfare tradeoffs. \\
Scope income-composition effects & Seat-pricing increments compare two optimized policies. HN/WC common-price baselines are certified; differentiated and JX designs have joint-search evidence. \\
Distinguish K objectives and availability & Commercial additions can enable profitable incumbent adjustments while lowering traveler surplus. JX welfare depends on schedule value, newly served periods and fleets. \\
Interpret the per-trip figure as a budget ratio & Profit sacrifice per restored L trip concerns broad fare and service restraint. A voucher does not reproduce benefits to other travelers. \\
Keep a participation safeguard & Lower realized-traveler cost ratios can coexist with vulnerable non-travel in the tested design library. \\
\bottomrule\end{tabular*}\par\medskip
\par\noindent\fontsize{8}{9.5}\selectfont Results depend on calibrated demand, costs and authority; empirical demand and full-route accounts are needed for implementation.
\end{minipage}\end{table}
